\section{Published post-foreclosure calculation}\label{sec:calculation}

We retain the published model and change no primitive. Following the circular-market structure of \citet{salop1979}, in each national market three firms are located one unit apart on a Salop circle of circumference three, consumers have density one, and transportation disutility is the square of shortest-arc distance. Each consumer buys one unit. Production cost is normalized to zero. When countries 1 and 2 form a standardization union, their firms can sell in either member country without conversion costs, while the outsider firm incurs conversion cost $c$ when selling there \citep[pp.~367--368, 373--374]{gandalshy2001}. Markets are segmented, so the price subgame can be studied one member market at a time.

Before foreclosure, the three-firm interior solution in a member market has
\begin{equation}\label{eq:interior-su}
 p_M=1+\frac{c}{5},\qquad p_O=1+\frac{3c}{5},\qquad
 q_M=1+\frac{c}{5},\qquad q_O=1-\frac{2c}{5}.
\end{equation}
Thus the outsider's interior share reaches zero at $c=5/2$. The correction below does not alter that threshold. It concerns the price game after the outsider has zero demand.

With firm 3 foreclosed, firms 1 and 2 face one another across both arcs joining their locations. One arc has length one and the other length two. On the short arc, measuring distance $x$ from firm 1 toward firm 2, the indifferent consumer solves
\[
 p_1+x^2=p_2+(1-x)^2,
\]
so
\begin{equation}\label{eq:shortarc}
 x_{12}^{S}=\frac12+\frac{p_2-p_1}{2}.
\end{equation}
On the long arc the corresponding condition is
\[
 p_1+x^2=p_2+(2-x)^2,
\]
which gives
\begin{equation}\label{eq:longarc}
 x_{12}^{L}=1+\frac{p_2-p_1}{4}.
\end{equation}
The long-arc price-difference coefficient is therefore one quarter, not one half. Appendix B of \citet[pp.~381--382]{gandalshy2001} instead uses the same one-half coefficient on the long arc as on the short arc before reporting the post-foreclosure profile $(3/2,3/2,c)$. The long-arc expression is inconsistent with the quadratic transportation primitive in the published article.\footnote{An earlier working-paper version uses linear transportation costs, under which the corresponding calculation belongs to a different specification. The published article uses quadratic transportation disutility; the comparison identifies the specification difference without attributing a cause to the published discrepancy. See \citet{gandalshy1996} and \citet[pp.~367--368]{gandalshy2001}.}

Whenever the outsider is irrelevant at both member boundaries and the two boundaries remain on the interior of their respective arcs, adding \eqref{eq:shortarc} and \eqref{eq:longarc} yields firm 1's two-member demand
\begin{equation}\label{eq:duodemand}
 q_1(p_1,p_2)=\frac32+\frac34(p_2-p_1),
\end{equation}
with the symmetric expression for firm 2. This formula is a demand-region formula, not a global demand formula for arbitrary price vectors; deviations that make the outsider relevant must be treated separately.

The published complete profile fails the Nash test throughout the strict post-foreclosure range. At $(p_1,p_2,p_3)=(3/2,3/2,c)$ with any $c>5/2$, let firm 1 raise its price to $3/2+\varepsilon$, where $\varepsilon>0$ is chosen small enough that the outsider remains foreclosed and $\varepsilon<1/2$. On the regular two-member branch, the profit gain is
\begin{align}\label{eq:published-profile-gain}
 \Delta\pi_1
 &=\left(\frac32+\varepsilon\right)
   \left(\frac32-\frac34\varepsilon\right)-\frac94 \\
 &=\frac38\varepsilon-\frac34\varepsilon^2
 =\frac38\varepsilon(1-2\varepsilon)>0.
\end{align}
Because $c-5/2>0$, such a deviation can be chosen within the foreclosure slack. Hence the complete Appendix-B profile $(3/2,3/2,c)$ is not a Nash equilibrium for the strict range $c>5/2$ considered here. At the exact checkpoint $c=4$, taking $\varepsilon=1/4$ raises profit from $9/4$ to $147/64$, a gain of $3/64$.

This correction must be stated at the level of the complete price profile. It does not imply that the member price $3/2$ can never be supported by another zero-sales outsider quote. The next section characterizes that distinction explicitly.
